% GENERATED FILE. Edit canonical lessons, then run npm run generate:books.
% canonical-source-hash: fnv1a64:00e7251680d3c5a3
% canonical-lessons: PA-C350-shauk, PA-C350-natak, PA-C350-filam, PA-C350-nachauna, PA-C350-parchauna, PA-R350-review-344-350

\chapter{Plays, Films and Fun}
\label{ch:pa-a2-plays-films-and-fun}
\hlchaptermodality{350}

\begin{chapteropening}
\textbf{By the end of this chapter:} \emph{I can talk about my hobbies and the shows I enjoy.}

Complete the last lesson of chapter 350: I can talk about my hobbies and the shows I enjoy.
\end{chapteropening}

\section[\texorpdfstring{\pa{ਸ਼ੌਕ} (shauk)}{ਸ਼ੌਕ (shauk)}]{\pa{ਸ਼ੌਕ} (shauk) — a hobby; a liking}

\label{lesson:PA-C350-shauk}



\begin{quote}
Before the new one: say the Punjabi for a poem, then the Punjabi for to praise.
\end{quote}

\subsection*{You'll want to know: \pa{ਸ਼ੌਕ}}
\textbf{\pa{ਸ਼ੌਕ}} — \emph{shauk} — \textquotedblleft{}a hobby; a liking\textquotedblright{}.

\begin{quote}
\pa{ਮੇਰਾ} \pa{ਸ਼ੌਕ} \pa{ਸੰਗੀਤ} \pa{ਹੈ।}
\end{quote}

\begin{quote}
\emph{merā shauk saṅgīt hai.}
\end{quote}

\begin{quote}
My hobby is music.
\end{quote}

A masculine noun.

\begin{cousinweb}
\textbf{\pa{ਸ਼ੌਕ}} (\emph{shauk}) is a loan, usually traced through Persian to Arabic \emph{shawq}, \textquotedblleft{}longing\textquotedblright{}.
\end{cousinweb}

\begin{grammarlens}[title={What you've built}]
One more word for this chapter.
\end{grammarlens}

\subsection*{Your turn}
\begin{itemize}
\raggedright
  \item \emph{Say it:} \emph{shauk}
  \item \emph{Say it:} \emph{shauk}, once more
  \item \emph{Say it:} \emph{merā shauk saṅgīt hai.} — My hobby is music.
  \item \emph{Recall:} say the Punjabi for a poem, then the Punjabi for to praise, then say \emph{shauk} again
\end{itemize}

\begin{tcolorbox}[breakable,colback=teal!4,colframe=teal!35!black,title={Before you move on}]
What does \pa{ਸ਼ੌਕ} mean? (\textquotedblleft{}A hobby; a liking\textquotedblright{}.) Say it once more.
\end{tcolorbox}
\section[\texorpdfstring{\pa{ਨਾਟਕ} (nāṭak)}{ਨਾਟਕ (nāṭak)}]{\pa{ਨਾਟਕ} (nāṭak) — a stage play; a drama}

\label{lesson:PA-C350-natak}



\begin{quote}
Before the new one: say the Punjabi for to praise, then the Punjabi for a hobby.
\end{quote}

\subsection*{You'll want to know: \pa{ਨਾਟਕ}}
\textbf{\pa{ਨਾਟਕ}} — \emph{nāṭak} — \textquotedblleft{}a stage play; a drama\textquotedblright{}.

\begin{quote}
\pa{ਅਸੀਂ} \pa{ਐਤਵਾਰ} \pa{ਨੂੰ} \pa{ਇੱਕ} \pa{ਨਾਟਕ} \pa{ਦੇਖਾਂਗੇ।}
\end{quote}

\begin{quote}
\emph{asīṁ aitvār nū̃ ikk nāṭak dekhāṅge.}
\end{quote}

\begin{quote}
We will watch a play on Sunday.
\end{quote}

A masculine noun.

\begin{cousinweb}
\textbf{\pa{ਨਾਟਕ}} (\emph{nāṭak}) is taken over from Sanskrit \emph{nāṭaka}, \textquotedblleft{}acting, dancing\textquotedblright{}, a tatsama (Monier-Williams).
\end{cousinweb}

\begin{grammarlens}[title={What you've built}]
One more word for this chapter.
\end{grammarlens}

\subsection*{Your turn}
\begin{itemize}
\raggedright
  \item \emph{Say it:} \emph{nāṭak}
  \item \emph{Say it:} \emph{nāṭak}, once more
  \item \emph{Say it:} \emph{asīṁ aitvār nū̃ ikk nāṭak dekhāṅge.} — We will watch a play on Sunday.
  \item \emph{Recall:} say the Punjabi for to praise, then the Punjabi for a hobby, then say \emph{nāṭak} again
  \item \emph{Write it:} copy \textbf{\pa{ਨਾਟਕ}} once, checking the whole word
\end{itemize}

\begin{tcolorbox}[breakable,colback=teal!4,colframe=teal!35!black,title={Before you move on}]
What does \pa{ਨਾਟਕ} mean? (\textquotedblleft{}A stage play; a drama\textquotedblright{}.) Say it once more.
\end{tcolorbox}
\section[\texorpdfstring{\pa{ਫ਼ਿਲਮ} (filam)}{ਫ਼ਿਲਮ (filam)}]{\pa{ਫ਼ਿਲਮ} (filam) — a film}

\label{lesson:PA-C350-filam}



\begin{quote}
Before the new one: say the Punjabi for a hobby, then the Punjabi for a stage play.
\end{quote}

\subsection*{You'll want to know: \pa{ਫ਼ਿਲਮ}}
\textbf{\pa{ਫ਼ਿਲਮ}} — \emph{filam} — \textquotedblleft{}a film\textquotedblright{}.

\begin{quote}
\pa{ਮੈਨੂੰ} \pa{ਪੁਰਾਣੀਆਂ} \pa{ਫ਼ਿਲਮਾਂ} \pa{ਪਸੰਦ} \pa{ਹਨ।}
\end{quote}

\begin{quote}
\emph{mainū̃ purāṇīāṁ filmāṁ pasand han.}
\end{quote}

\begin{quote}
I like old films.
\end{quote}

A feminine noun.

\begin{cousinweb}
\textbf{\pa{ਫ਼ਿਲਮ}} (\emph{filam}) is borrowed from English \emph{film}.
\end{cousinweb}

\begin{grammarlens}[title={What you've built}]
One more word for this chapter.
\end{grammarlens}

\subsection*{Your turn}
\begin{itemize}
\raggedright
  \item \emph{Say it:} \emph{filam}
  \item \emph{Say it:} \emph{filam}, once more
  \item \emph{Say it:} \emph{mainū̃ purāṇīāṁ filmāṁ pasand han.} — I like old films.
  \item \emph{Recall:} say the Punjabi for a hobby, then the Punjabi for a stage play, then say \emph{filam} again
  \item \emph{Write it:} copy \textbf{\pa{ਫ਼ਿਲਮ}} once, checking the whole word
\end{itemize}

\begin{tcolorbox}[breakable,colback=teal!4,colframe=teal!35!black,title={Before you move on}]
What does \pa{ਫ਼ਿਲਮ} mean? (\textquotedblleft{}A film\textquotedblright{}.) Say it once more.
\end{tcolorbox}
\section[\texorpdfstring{\pa{ਨਚਾਉਣਾ} (nachāuṇā)}{ਨਚਾਉਣਾ (nachāuṇā)}]{\pa{ਨਚਾਉਣਾ} (nachāuṇā) — to make someone dance; to get someone dancing}

\label{lesson:PA-C350-nachauna}



\begin{quote}
Before the new one: say the Punjabi for a stage play, then the Punjabi for a film.
\end{quote}

\subsection*{You'll want to know: \pa{ਨਚਾਉਣਾ}}
\textbf{\pa{ਨਚਾਉਣਾ}} — \emph{nachāuṇā} — \textquotedblleft{}to make someone dance; to get someone dancing\textquotedblright{}.

\begin{quote}
\pa{ਢੋਲ} \pa{ਦੀ} \pa{ਥਾਪ} \pa{ਸਭ} \pa{ਨੂੰ} \pa{ਨਚਾਉਂਦੀ} \pa{ਹੈ।}
\end{quote}

\begin{quote}
\emph{ḍhol dī thāp sabh nū̃ nachāundī hai.}
\end{quote}

\begin{quote}
The beat of the drum gets everyone dancing.
\end{quote}

Past (I …): \pa{ਮੈਂ} \pa{ਬੱਚਿਆਂ} \pa{ਨੂੰ} \pa{ਨਚਾਇਆ}, \emph{maiṁ bachchiāṁ nū̃ nachāiā}.

\begin{cousinweb}
\textbf{\pa{ਨਚਾਉਣਾ}} (\emph{nachāuṇā}) is the causative partner of \emph{naccṇā}, \textquotedblleft{}to dance\textquotedblright{}, which goes back to Sanskrit \emph{nṛtyati}, \textquotedblleft{}dances\textquotedblright{}; CDIAL lists Punjabi \emph{naccṇā} under it (CDIAL 7583).
\end{cousinweb}

\begin{grammarlens}[title={What you've built}]
One more word for this chapter.
\end{grammarlens}

\subsection*{Your turn}
\begin{itemize}
\raggedright
  \item \emph{Say it:} \emph{nachāuṇā}
  \item \emph{Say it:} \emph{nachāuṇā}, once more
  \item \emph{Say it:} \emph{ḍhol dī thāp sabh nū̃ nachāundī hai.} — The beat of the drum gets everyone dancing.
  \item \emph{Recall:} say the Punjabi for a stage play, then the Punjabi for a film, then say \emph{nachāuṇā} again
  \item \emph{Write it:} copy \textbf{\pa{ਨਚਾਉਣਾ}} once, checking the whole word
\end{itemize}

\begin{tcolorbox}[breakable,colback=teal!4,colframe=teal!35!black,title={Before you move on}]
What does \pa{ਨਚਾਉਣਾ} mean? (\textquotedblleft{}To make someone dance; to get someone dancing\textquotedblright{}.) Say it once more.
\end{tcolorbox}
\section[\texorpdfstring{\pa{ਪਰਚਾਉਣਾ} (parchāuṇā)}{ਪਰਚਾਉਣਾ (parchāuṇā)}]{\pa{ਪਰਚਾਉਣਾ} (parchāuṇā) — to entertain; to amuse}

\label{lesson:PA-C350-parchauna}



\begin{quote}
Before the new one: say the Punjabi for a film, then the Punjabi for to make someone dance.
\end{quote}

\subsection*{You'll want to know: \pa{ਪਰਚਾਉਣਾ}}
\textbf{\pa{ਪਰਚਾਉਣਾ}} — \emph{parchāuṇā} — \textquotedblleft{}to entertain; to amuse\textquotedblright{}.

\begin{quote}
\pa{ਨਾਨੀ} \pa{ਬੱਚਿਆਂ} \pa{ਨੂੰ} \pa{ਗੀਤਾਂ} \pa{ਨਾਲ} \pa{ਪਰਚਾਉਂਦੀ} \pa{ਹੈ।}
\end{quote}

\begin{quote}
\emph{nānī bachchiāṁ nū̃ gītāṁ nāl parchāundī hai.}
\end{quote}

\begin{quote}
Grandmother keeps the children amused with songs.
\end{quote}

Past (I …): \pa{ਮੈਂ} \pa{ਮਹਿਮਾਨਾਂ} \pa{ਨੂੰ} \pa{ਪਰਚਾਇਆ}, \emph{maiṁ mahimānāṁ nū̃ parchāiā}.

\begin{cousinweb}
\textbf{\pa{ਪਰਚਾਉਣਾ}} (\emph{parchāuṇā}) is inherited from a reconstructed \emph{*prarajyate}, \textquotedblleft{}is pleased, is excited\textquotedblright{}; CDIAL lists Punjabi \emph{parcāuṇā}, \textquotedblleft{}to divert, entertain\textquotedblright{}, under it (CDIAL 8737).
\end{cousinweb}

\begin{grammarlens}[title={What you've built}]
One more word for this chapter.
\end{grammarlens}

\subsection*{Your turn}
\begin{itemize}
\raggedright
  \item \emph{Say it:} \emph{parchāuṇā}
  \item \emph{Say it:} \emph{parchāuṇā}, once more
  \item \emph{Say it:} \emph{nānī bachchiāṁ nū̃ gītāṁ nāl parchāundī hai.} — Grandmother keeps the children amused with songs.
  \item \emph{Recall:} say the Punjabi for a film, then the Punjabi for to make someone dance, then say \emph{parchāuṇā} again
  \item \emph{Write it:} copy \textbf{\pa{ਪਰਚਾਉਣਾ}} once, checking the whole word
\end{itemize}

\begin{tcolorbox}[breakable,colback=teal!4,colframe=teal!35!black,title={Before you move on}]
What does \pa{ਪਰਚਾਉਣਾ} mean? (\textquotedblleft{}To entertain; to amuse\textquotedblright{}.) Say it once more.
\end{tcolorbox}
\section[(dialogue)]{Looking back: from We Don't Agree to Plays, Films and Fun}

\label{lesson:PA-R350-review-344-350}



\begin{quote}
Say the Punjabi for an argument and a debate.
\end{quote}

\subsection*{Your turn}
Say these aloud:

\begin{itemize}
\raggedright
  \item an argument, a debate, a dispute, stubbornness, an objection
  \item a condition, a decision, an opinion, to settle a matter, a proposal
  \item a fault, to hide something, an accusation, helpless, a pity
  \item a drum, to play an instrument, rhythm, a drumbeat, a clap
  \item a musical instrument, a tabla, a tune, music practice, to echo
  \item a singer, a flute, a harmonium, a poem, to praise
  \item a hobby, a stage play, a film, to make someone dance, to entertain
\end{itemize}

\begin{tcolorbox}[breakable,colback=teal!4,colframe=teal!35!black,title={Before you move on}]
An argument? (\textbf{\pa{ਦਲੀਲ}}, \emph{dalīl}.) A debate? (\textbf{\pa{ਬਹਿਸ}}, \emph{bahis}.) A dispute? (\textbf{\pa{ਤਕਰਾਰ}}, \emph{takrār}.) Stubbornness? (\textbf{\pa{ਜ਼ਿੱਦ}}, \emph{zidd}.) An objection? (\textbf{\pa{ਇਤਰਾਜ਼}}, \emph{itrāz}.) A condition? (\textbf{\pa{ਸ਼ਰਤ}}, \emph{sharat}.) A decision? (\textbf{\pa{ਫ਼ੈਸਲਾ}}, \emph{faislā}.) An opinion? (\textbf{\pa{ਰਾਏ}}, \emph{rāe}.) To settle a matter? (\textbf{\pa{ਮੁਕਾਉਣਾ}}, \emph{mukāuṇā}.) A proposal? (\textbf{\pa{ਤਜਵੀਜ਼}}, \emph{tajvīz}.) A fault? (\textbf{\pa{ਕਸੂਰ}}, \emph{kasūr}.) To hide something? (\textbf{\pa{ਲੁਕਾਉਣਾ}}, \emph{lukāuṇā}.) An accusation? (\textbf{\pa{ਇਲਜ਼ਾਮ}}, \emph{ilzām}.) Helpless? (\textbf{\pa{ਮਜਬੂਰ}}, \emph{majbūr}.) A pity? (\textbf{\pa{ਅਫ਼ਸੋਸ}}, \emph{afsos}.) A drum? (\textbf{\pa{ਢੋਲ}}, \emph{ḍhol}.) To play an instrument? (\textbf{\pa{ਵਜਾਉਣਾ}}, \emph{vajāuṇā}.) Rhythm? (\textbf{\pa{ਤਾਲ}}, \emph{tāl}.) A drumbeat? (\textbf{\pa{ਥਾਪ}}, \emph{thāp}.) A clap? (\textbf{\pa{ਤਾੜੀ}}, \emph{tāṛī}.) A musical instrument? (\textbf{\pa{ਸਾਜ਼}}, \emph{sāz}.) A tabla? (\textbf{\pa{ਤਬਲਾ}}, \emph{tablā}.) A tune? (\textbf{\pa{ਧੁਨ}}, \emph{dhun}.) Music practice? (\textbf{\pa{ਰਿਆਜ਼}}, \emph{riāz}.) To echo? (\textbf{\pa{ਗੂੰਜਣਾ}}, \emph{gūnjṇā}.) A singer? (\textbf{\pa{ਗਾਇਕ}}, \emph{gāik}.) A flute? (\textbf{\pa{ਵੰਝਲੀ}}, \emph{vanjhlī}.) A harmonium? (\textbf{\pa{ਵਾਜਾ}}, \emph{vājā}.) A poem? (\textbf{\pa{ਕਵਿਤਾ}}, \emph{kavitā}.) To praise? (\textbf{\pa{ਸਰਾਹੁਣਾ}}, \emph{sarāhuṇā}.) A hobby? (\textbf{\pa{ਸ਼ੌਕ}}, \emph{shauk}.) A stage play? (\textbf{\pa{ਨਾਟਕ}}, \emph{nāṭak}.) A film? (\textbf{\pa{ਫ਼ਿਲਮ}}, \emph{filam}.) To make someone dance? (\textbf{\pa{ਨਚਾਉਣਾ}}, \emph{nachāuṇā}.) To entertain? (\textbf{\pa{ਪਰਚਾਉਣਾ}}, \emph{parchāuṇā}.)
\end{tcolorbox}
